%==============================================================================
% VALIDATION - PDB Structure Comparison
%==============================================================================
\subsection{PDB Structure Validation}
\label{sec:pdb_validation}

To provide independent validation beyond our benchmark dataset, we compared GlycoSMILES2BAP output against known correct structures from the Protein Data Bank (PDB). This validation approach addresses the limitation of AF3 model access while leveraging experimentally verified glycan structures.

\subsubsection{Reference Structure Selection}

We selected four PDB structures representing critical stereochemistry test cases:

\begin{table}[h]
\centering
\caption{PDB Reference Structures for Validation}
\begin{tabular}{llll}
\toprule
PDB ID & Glycan Type & Key Validation Point & Reference \\
\midrule
5NSC & Fucosylated & FUC = alpha-L-fucose & \citep{frenz2018} \\
1L6R & Lactose & GAL C4 axial hydroxyl & \citep{lactose_ref} \\
2VXR & Sialyllactose & SIA anomeric = C2 & \citep{sialic_ref} \\
5K65 & M3 N-glycan & Branch topology & \citep{jo2011} \\
\bottomrule
\end{tabular}
\label{tab:pdb_refs}
\end{table}

\subsubsection{Validation Method}

For each PDB structure, we:
\begin{enumerate}
\item Extracted the glycan IUPAC notation from literature
\item Generated CCD codes using our CCD Mapper module
\item Compared anomeric carbon and ring oxygen positions
\item Verified configuration (D/L) for each monosaccharide
\end{enumerate}

\subsubsection{Validation Results}

All four test cases passed validation (100\% accuracy):

\begin{table}[h]
\centering
\caption{PDB Validation Results}
\begin{tabular}{llcccc}
\toprule
Test Case & Expected CCD & Actual CCD & Anomeric & Ring O & Status \\
\midrule
Fucosylated (5NSC) & FUC & FUC & C1 & O5 & \checkmark \\
Lactose (1L6R) & GAL & GAL & C1 & O5 & \checkmark \\
Sialyllactose (2VXR) & SIA & SIA & C2 & O6 & \checkmark \\
M3 N-glycan (5K65) & MAN/BMA/NAG & MAN/BMA/NAG & C1 & O5 & \checkmark \\
\bottomrule
\end{tabular}
\label{tab:pdb_results}
\end{table}

\textbf{Critical Finding: Sialic Acid C2 Anomeric Position}

The most important validation was the sialyllactose structure (PDB: 2VXR). Sialic acids (Neu5Ac) are ketoses with the anomeric carbon at C2, not C1 as in aldoses. The SMILES format typically fails this distinction, placing the glycosidic bond at C1 incorrectly.

GlycoSMILES2BAP correctly:
\begin{itemize}
\item Identified SIA anomeric carbon as C2 (ketose property)
\item Set ring oxygen to O6 (not O5)
\item Generated correct CCD code (SIA for alpha-D-Neu5Ac)
\end{itemize}

This validates the core algorithm for handling ketose monosaccharides, addressing the stereochemistry error identified by Huang et al.\ (2025).

%==============================================================================
% RESULTS - Benchmark
%==============================================================================
\section{Results}

\subsection{Benchmark Dataset}

We constructed a benchmark dataset of 50 glycan structures covering diverse categories:

\begin{table}[h]
\centering
\begin{tabular}{lcl}
\toprule
Category & Count & Examples \\
\midrule
Linear glycans & 15 & LNnT, LNT, maltose polymers \\
N-glycans & 20 & M3-M9, G0-G2, fucosylated \\
O-glycans & 10 & Tn, T, sialyl-T antigens \\
Complex & 5 & Sialylated, branched, rare sugars \\
\bottomrule
\end{tabular}
\caption{Benchmark dataset composition}
\end{table}

\subsection{Benchmark Performance}

\begin{table}[h]
\centering
\caption{Benchmark Performance Comparison}
\begin{tabular}{lccccc}
\toprule
Metric & Ours & 95\% CI &